\chapter{Getting Zeno}\label{ch-getting-zeno}

Zeno is one program, called \verb|zeno|. It checks your programs, runs
them, compiles them, tests them, lays them out, and talks to your
editor. This chapter puts it on your machine.

There are two ways to get it: download a \verb|zeno| that is already
built, or build it yourself from its source. The first takes a minute
and needs nothing installed. The second needs Rust, and gives you the
compiler's source to read beside \hyperref[ch-inside]{Inside Zeno}.

\section{Downloading it}

Every release of Zeno is built for ten kinds of machine, and put on
the releases page of the repository:

\begin{center}
\url{https://github.com/partham20/zeno-lang/releases}
\end{center}

Pick the file for your machine. The number after \verb|zeno-| is the
version; here it is written \verb|V|.

\begin{center}
\begin{tabular}{P{6.6cm}l}
\toprule
Machine & File \\
\midrule
Linux on most PCs and servers (x86-64) & \verb|zeno-V-linux-x86_64.tar.gz| \\
Linux on 32-bit x86 & \verb|zeno-V-linux-x86.tar.gz| \\
Raspberry Pi 3, 4 or 5 with a 64-bit system, and ARM servers & \verb|zeno-V-linux-arm64.tar.gz| \\
Raspberry Pi 2 or later with a 32-bit system (ARMv7) & \verb|zeno-V-linux-armv7.tar.gz| \\
Raspberry Pi 1 or Zero (ARMv6) & \verb|zeno-V-linux-armv6.tar.gz| \\
Mac with Apple silicon (M1 and later) & \verb|zeno-V-macos-arm64.tar.gz| \\
Mac with an Intel processor & \verb|zeno-V-macos-x86_64.tar.gz| \\
Windows on most PCs (x86-64) & \verb|zeno-V-windows-x86_64.zip| \\
Windows on 32-bit x86 & \verb|zeno-V-windows-x86.zip| \\
Windows on ARM & \verb|zeno-V-windows-arm64.zip| \\
\bottomrule
\end{tabular}
\end{center}

Inside is one folder with \verb|zeno| in it (\verb|zeno.exe| on
Windows), and, from version 0.2, \verb|zeno.wasm| beside it, which
\hyperref[ch-web]{A Program in a Web Page} needs. \verb|zeno| is one
file that needs nothing else installed: no Rust, no libraries.

\subsection{On Linux and the Raspberry Pi}

Unpack the file, and copy \verb|zeno| into a folder that is on your
\verb|PATH|. \verb|~/.local/bin| is one on most systems, and needs no
\verb|sudo|:

\begin{zenofrag}[lang=shell]
$ tar -xzf zeno-0.1.0-linux-x86_64.tar.gz
$ mkdir -p ~/.local/bin
$ cp zeno-0.1.0-linux-x86_64/zeno* ~/.local/bin/
$ zeno --version
zeno 0.1.0
\end{zenofrag}

If the shell says \verb|zeno: command not found|, that folder is not on
your \verb|PATH| yet. Add the line \verb|export PATH="$HOME/.local/bin:$PATH"|
to \verb|~/.bashrc|, and open a new terminal.

The Linux builds are linked with everything they need inside them, so
the same file runs on Debian, Ubuntu, Fedora, Alpine or Raspberry Pi
OS, old or new.

\subsection{On macOS}

The same steps work in the Terminal, with \verb|/usr/local/bin| as the
folder:

\begin{zenofrag}[lang=shell]
$ tar -xzf zeno-0.1.0-macos-arm64.tar.gz
$ sudo cp zeno-0.1.0-macos-arm64/zeno* /usr/local/bin/
$ xattr -d com.apple.quarantine /usr/local/bin/zeno
$ zeno --version
\end{zenofrag}

The \verb|xattr| line is needed once. A program downloaded with a
browser is held back by macOS until Apple has signed it, and Zeno is
not signed; the line tells macOS that you trust this one.

\subsection{On Windows}

Unzip the file: right-click it in Explorer and choose
\emph{Extract All}. Make a folder for programs of your own, such as
\verb|C:\Tools|, and move \verb|zeno.exe| (and \verb|zeno.wasm|) into
it. Then put that folder on the \verb|PATH|: open the Start menu, type
\emph{environment variables}, choose \emph{Edit environment variables
for your account}, select \verb|Path|, press \emph{Edit} and
\emph{New}, and type the folder. A new Command Prompt or PowerShell
window then finds it:

\begin{zenofrag}[lang=shell]
> zeno --version
zeno 0.1.0
\end{zenofrag}

Windows may warn that the program is from an unknown publisher, for
the same reason as on macOS. \emph{More info}, then \emph{Run anyway},
lets it start.

\subsection{Checking the download}

Beside the files on the releases page is \verb|SHA256SUMS.txt|, with
the fingerprint of each one. To be sure that the file you have is the
one that was built:

\begin{zenofrag}[lang=shell]
$ sha256sum --check --ignore-missing SHA256SUMS.txt
zeno-0.1.0-linux-x86_64.tar.gz: OK
\end{zenofrag}

On macOS the command is \verb|shasum -a 256 -c|, and in PowerShell
\verb|Get-FileHash| prints a file's fingerprint to compare by eye.

\section{What else you may need}

\verb|zeno run|, \verb|check|, \verb|test|, \verb|format| and the
language server need nothing more. Three commands hand work to
programs of other people:

\begin{itemize}
\item \verb|zeno build| writes C and gives it to a C compiler: \verb|cc|,
  or the one the variable \verb|CC| names. On Linux, install
  \verb|gcc| or \verb|clang| from the system's packages; on macOS,
  \verb|xcode-select --install|; on Windows, MinGW's \verb|gcc|, with
  \verb|set CC=gcc| (\hyperref[ch-compiling]{Compiling}).
\item \verb|zeno build --llvm| needs \verb|clang|.
\item \verb|zeno fetch| needs \verb|git|, to get packages that live in
  other people's repositories (\hyperref[ch-packages]{Packages}).
\end{itemize}

\section{Building it from source}

To build Zeno yourself you need:

\begin{itemize}
\item \textbf{Rust}, with its build tool \verb|cargo|. The usual way
  to get both is \verb|rustup|, from \url{https://rustup.rs}. Zeno
  uses no libraries beyond Rust itself, so nothing else is
  downloaded when it is built.
\item \textbf{git}, to get the source.
\end{itemize}

Zeno is built and tested on Linux, macOS and Windows, with the same
steps on each. Get the source and build it:

\begin{zenofrag}[lang=shell]
git clone https://github.com/partham20/zeno-lang
cd zeno-lang
cargo build --release
\end{zenofrag}

The finished program is \verb|target/release/zeno|. To use it from
any folder, either put that folder on your \verb|PATH|, or let cargo
copy the program into its own folder of programs, which \verb|rustup|
has already put on your \verb|PATH|:

\begin{zenofrag}[lang=shell]
cargo install --path .
\end{zenofrag}

For web pages, build \verb|zeno.wasm| as well. It is the same
compiler, built for a browser:

\begin{zenofrag}[lang=shell]
rustup target add wasm32-unknown-unknown
cargo build --release --target wasm32-unknown-unknown
\end{zenofrag}

Check that it works:

\begin{session}
$ zeno --version
\end{session}

If you want to be sure that the compiler you built is a good one, run
its own tests. There are several hundred of them, and each program in
them is run by every engine and built every way Zeno can build it:

\begin{zenofrag}[lang=shell]
cargo test --release
\end{zenofrag}

\section{Your first project}

\verb|zeno new| makes a folder with a small program in it, ready to
run:

\begin{session}
$ zeno new greeter
$ zeno run greeter/main.zeno
$ zeno test greeter
\end{session}

Three files were made. \verb|main.zeno| is where the program starts.
\verb|greeting.zeno| is a \emph{module}, a second file of the program,
with a test at its end. \verb|zeno.package| names the program and,
later, the packages it needs. You will meet each of them properly:
modules in \hyperref[ch-modules]{Programs in Several Files}, packages
in \hyperref[ch-packages]{Packages}, and tests in
\hyperref[ch-testing]{Testing}. For now, a program can be a single
file with any name.

\section{The commands}

Everything Zeno does is a word after \verb|zeno|. These are the ones
you will use most; \hyperref[app-commands]{The zeno Command} has all
of them with every option.

\begin{center}
\begin{tabular}{lP{9cm}}
\toprule
Command & What it does \\
\midrule
\verb|zeno run file.zeno| & Check the program, and run it if there is nothing wrong \\
\verb|zeno check file.zeno| & Only check it, and report every mistake \\
\verb|zeno test| & Run the tests written in the files of a folder \\
\verb|zeno build file.zeno| & Compile it to C, and then to a program of its own \\
\verb|zeno format file.zeno| & Lay the file out the one way Zeno lays programs out \\
\verb|zeno new name| & Make a new program in a new folder \\
\verb|zeno fetch| & Get the packages a program needs \\
\verb|zeno lsp| & The language server that an editor starts \\
\bottomrule
\end{tabular}
\end{center}

\section{An editor}

Any editor that can save plain text will do. Zeno files end in
\verb|.zeno| and are UTF-8 text. The Zeno repository has colours and
setup for VS Code, Vim, Neovim and Helix in its \verb|editors| folder,
and \verb|zeno lsp| lets those editors show mistakes as you type, jump
to where a name was declared, show a name's type, and lay the file out
when you save it. \hyperref[ch-tools]{The Tools} shows how to set each
of them up.

\begin{tryit}
\begin{enumerate}
\item Run \verb|zeno --help| and read the list of commands. Which
  ones need a C compiler?
\item Change the text that \verb|greeting.zeno| gives back, and run
  \verb|zeno test greeter| again. Read what a failed test prints, then
  change the test so that it passes.
\item Run \verb|zeno new 2fast| and \verb|zeno new while|. What does
  Zeno say, and what does it suggest instead?
\end{enumerate}
\end{tryit}
